\section{Current Numerical Validation}

The repository currently contains automated tests and numerical examples that verify the solver independently of future experiments. These results establish that the computational core is functioning before experimental validation is attempted.

\subsection{Automated Tests}

The current test suite passes all available tests:
\begin{equation*}
\texttt{58 passed in 45 seconds}.
\end{equation*}
The tests cover grid indexing, sparse operator assembly, axis regularity stencils, boundary conditions, Poisson sign convention, electric-field reconstruction, interpolation, curvature, residual behavior, Taylor-angle benchmark, space-charge closures, optimization sanity checks, app-backend outputs, immersed boundary method accuracy, and Taylor onset verification.

\subsection{Classical Taylor-Angle Benchmark}

The analytical benchmark routine computes
\begin{equation*}
\alpha_T = \SI{49.2900892921}{\degree},
\end{equation*}
matching the classical Taylor half-angle. This confirms that the theoretical benchmark is correctly represented in code.

\subsection{Manufactured Poisson Convergence}

A manufactured solution is used to verify the Poisson solver:
\begin{equation}
\ph(r,z)=r^4+z^4+r^2z^2.
\end{equation}
For this function,
\begin{equation}
\nabla^2\ph=18r^2+16z^2.
\end{equation}
The existing convergence data are shown in Table~\ref{tab:poisson-convergence}.

\begin{table}[H]
\centering
\caption{Manufactured Poisson convergence results from the current solver report.}
\label{tab:poisson-convergence}
\begin{tabular}{rrrrr}
\toprule
Grid & L2 error & Linf error & Nonzeros & Solve time (s) \\
\midrule
$25\times 27$ & $8.104\times 10^{-4}$ & $1.353\times 10^{-3}$ & 3050 & 0.0010 \\
$49\times 53$ & $1.992\times 10^{-4}$ & $3.389\times 10^{-4}$ & 12338 & 0.0027 \\
$97\times 105$ & $4.936\times 10^{-5}$ & $8.479\times 10^{-5}$ & 49634 & 0.0134 \\
\bottomrule
\end{tabular}
\end{table}

The decreasing error under refinement supports correctness of the finite-difference discretization.

\subsection{Voltage Scaling}

In a parallel-plate Laplace case, doubling voltage from \SI{500}{V} to \SI{1000}{V} produces an electric-field ratio of approximately 2 and an implied Maxwell-pressure ratio of approximately 4. This agrees with $\Emag\propto V$ and $p_E\propto E^2$.

\subsection{Residual Sanity Check}

For a cylindrical interface with zero electric field, capillary pressure is constant and is removed by the mean pressure offset. The current numerical report gives an RMS residual of approximately
\begin{equation*}
8.58\times 10^{-16}\ \si{Pa},
\end{equation*}
which is effectively numerical zero.

\subsection{Immersed Boundary Method (V3 Milestone)}

The solver now implements a sharp immersed boundary method with fractional-distance Dirichlet stencils for the conductor-gas interface. Key achievements include:

\paragraph{Second-Order Convergence.} Richardson extrapolation on a smooth manufactured problem demonstrates that the immersed Laplace operator converges at approximately second order in grid spacing, measured L2 order $\approx 2.0$ over three refinement levels. This replaces the previous staircase boundary representation, which did not converge cleanly.

\paragraph{Onset Voltage Projection.} The optimization objective now projects out the onset voltage amplitude per candidate shape rather than using a fixed applied voltage. For a candidate cone, the balance voltage $V_0^*$ is determined by minimizing the weighted Young--Laplace--Maxwell residual $\mathcal{R} = a - u^* b$, where $u^* = \langle ab \rangle_w / \langle b^2 \rangle_w$ and $V_0^* = \sqrt{2u^*/\varepsilon_0}$. This eliminates the degeneracy observed in the fixed-voltage objective.

\paragraph{Grounded-Box Free Boundary.} With onset projection, the immersed optimizer on a grounded rectangular domain converges to a half-angle of approximately \SI{44}{\degree} with onset voltage $V_0^* \approx \SI{28.9}{kV}$ for the meter-scale demonstration geometry. The angle direction shows a clean V-shaped landscape with an interior minimum.

\paragraph{Imposed-Taylor Identity Verification.} The Taylor amplitude identity
\begin{equation}
\frac{1}{2}\varepsilon_0 V_0^2 = \int_{\text{interface}} \gamma \kappa \, \mathrm{d}A
\end{equation}
is verified by imposing the exact analytic Taylor potential on the domain boundary with the rounded cone as the zero equipotential. The measured amplitude ratio is $1.009$ (within 1\% of unity), and the projected-residual argmin occurs at \SI{50.0}{\degree}, approximately \SI{0.7}{\degree} from the ideal \SI{49.29}{\degree}. This small systematic offset is attributed to the cap-flank geometry and discretization floor, and is stable under grid refinement. The strong result is the amplitude identity itself.

\paragraph{Normal-Field Reconstruction.} The solver uses a cubic-exact one-sided stencil to reconstruct the normal electric field $E_n = -(\partial \phi / \partial n)$ at the immersed boundary:
\begin{equation}
E_n = -\frac{-11V_b + 18V(d) - 9V(2d) + 2V(3d)}{6d},
\end{equation}
where $V_b$ is the boundary potential, $V(d)$, $V(2d)$, $V(3d)$ are full-cell samples along the outward normal, and $d = \max(\Delta r, \Delta z)$. Accuracy tests show $E_n$ error $\leq 1\%$ on refined grids for smooth profiles.

These results establish that the immersed free-boundary solver is verified and ready for comparison against experimental cone profiles and onset voltages.
